\begin{phasebox}
\textbf{هدف:} دریافت پیام برد و ذخیره آن در پایگاه داده.\\
\textbf{پیش‌نیاز:} قرارداد پیام و توابع پایگاه داده آماده باشند.\\
\textbf{معیار پایان:} پیام معتبر ذخیره و پیام نامعتبر رد شود.
\end{phasebox}

\section{محل کد}
کد دریافت پیام‌ها و منطق اصلی برنامه در پوشه
\lr{\texttt{code/backend}}
قرار می‌گیرد. این بخش همان بک‌اند پروژه است؛ یعنی برنامه پایتونی که بین دستگاه، پایگاه داده و برنامه وب قرار دارد.

\section{وظایف بک‌اند}
\begin{enumerate}
    \item اتصال به کارگزار
    \lr{MQTT}
    \item دریافت پیام داده یا وضعیت؛
    \item تبدیل پیام
    \lr{JSON}
    و بررسی فیلدهای آن؛
    \item فراخوانی تابع مناسب از پوشه پایگاه داده؛
    \item ثبت خطا و وضعیت اجرا.
\end{enumerate}

\section{قاعده مهم}
کد دریافت پیام نباید دستور مستقیم
\lr{SQL}
داشته باشد. این کد فقط تابع ذخیره‌سازی تعریف‌شده در پوشه پایگاه داده را فراخوانی می‌کند.

\section{آزمایش‌ها}
ابتدا پیام ثابت بدون برد ارسال می‌شود. پیام معتبر باید ذخیره شود و پیام دارای فیلد ناقص نباید وارد پایگاه داده شود. پس از موفقیت این تست، برد جای فرستنده خط فرمان قرار می‌گیرد.
